\section{Overall Prototype Feasibility}
\label{sec:ch4-overall-feasibility}

The combined results show that the TinyML Smart Plug achieved the intended
hybrid behavior under controlled laboratory conditions. The calibrated voltage,
current, and temperature measurements were accurate enough to support the
implemented thresholds. The overload tests showed warning behavior followed by
relay trip only when sustained overload persisted. The socket-heating tests
showed that degraded contact can cause faster and more severe heating than
current magnitude alone suggests. The arc-fault tests showed that waveform
features and Random Forest inference can support local relay action for
series-arc events in the tested load set.

The prototype is best interpreted as complementary protection. It adds
socket-level monitoring of overload persistence, thermal rise, and waveform
disturbance, but it is not a replacement for circuit breakers, AFCIs, GFCIs,
proper wiring practice, contact inspection, or certified protection devices.
This boundary matters because the tests were conducted on a laboratory bench,
the thermal estimator was deliberately conservative, and the TinyML model was
trained and evaluated on a finite set of loads and arc conditions.

The results nevertheless support the study objective: a compact
microcontroller-based smart plug can combine calibrated sensing, deterministic
threshold logic, compensated socket-temperature estimation, local TinyML
inference, relay latching, and PWA monitoring in a single prototype. The next
engineering step is not simply to lower thresholds or add more model
complexity; it is to broaden the dataset, validate the thermal estimator
against additional socket constructions, test across more appliance families,
and compare behavior against applicable protection-device standards.
